\documentclass[
10pt, % Main document font size
a4paper, % Paper type, use 'letterpaper' for US Letter paper
oneside, % One page layout (no page indentation)
%twoside, % Two page layout (page indentation for binding and different headers)
headinclude,footinclude, % Extra spacing for the header and footer
BCOR5mm, % Binding correction
]{scrartcl}
\input{stru/structure.tex}
\hyphenation{Fortran hy-phen-ation} 

\title{\normalfont{FTML practical session 6}} 
\author{\spacedlowsmallcaps{}
} 
\date{\today}
\begin{document}
\renewcommand{\sectionmark}[1]{\markright{\spacedlowsmallcaps{#1}}}
\lehead{\mbox{\llap{\small\thepage\kern1em\color{halfgray} \vline}\color{halfgray}\hspace{0.5em}\rightmark\hfil}} 
\pagestyle{scrheadings}
\maketitle 
\setcounter{tocdepth}{2} 

\begin{figure}[htpb]
    \centering
    \includegraphics[width=1\linewidth]{Gradient_descent}
    \label{fig:risks}
\end{figure}

\tableofcontents

\newpage


\section{\large\color{Blue}Gradient descent and line search for a least squares problem}

We come back to a least squares problem, that we will use to show an example of gradient descent algorithm, with a non-constant learning rate.

\subsection{Setting}

\subsubsection{Notations}%
\label{ssub:Notations}


We recall the setting and the notations here:

\begin{itemize}
    \item $ \mathcal{X}  = \mathbb{R}^d$
    \item $ \mathcal{Y} = \mathbb{R} $
    \item $ \theta = \mathbb{R}^d $ defines the estimator, which is
        the function $x\rightarrow x^T\theta$.
    \item design matrix: $X\in \mathbb{R}^{n,d}$
    \item vector of outputs (labels): $y \in \mathbb{R}^n$.
\end{itemize}

We want to minimize the function $f$ representing the empirical risk:

\begin{equation}
    f( \theta) = \frac{1}{2n} || X \theta - y||^2
\end{equation}

The $ \frac{1}{2}$ is just convenient here because it disappears in the gradient.

\subsubsection{Gradient descent}

We recall that the gradient of $f$ and its Hessian write:

    \begin{equation}
	\begin{aligned}
	    \label{eq:}
        \nabla_{\theta}f(\theta) &=  \frac{1}{n} X^T(X\theta-y)\\
	    &= H \theta - \frac{1}{n} X^Ty
	\end{aligned}
    \end{equation}

    \begin{equation}
        H = \frac{1}{n} X^TX
    \end{equation}

An \textbf{iteration} of the gradient algorithm writes:
\begin{equation}
    \theta_{t+1} = \theta_{t}- \gamma \nabla_{\theta}f(\theta_t)
\end{equation}

where $\gamma>0$ is the learning rate.

\subsubsection{Global minimizers}

We consider the global minimizers of $f$, noted $\theta^*$. They all necessary verify

    \begin{equation}
        \nabla_{\theta}f(\theta^*) = 0
    \end{equation}

    Which means that 

    \begin{equation}
        H\theta* = \frac{1}{n} X^Ty
    \end{equation}

    As $\theta\rightarrow f(\theta)$ is convex, we know that this
    condition is also sufficient: any vector that cancels the gradient is a global
    minimum.

If $H$ is not invertible, there might be several minimizers. However, if $f$ is
strongly convex, then $H$ is invertible and $ \theta^*$ is unique, we note it $\theta^*$. We will
consider such a case here. In general, $H$ is a positive semi-definite matrix,
and here we hence consider a case where it is definite positive.

\subsection{Vanilla gradient descent}

{\color{Purple}Implement a gradient descent (GD) on the data contained in \textbf{{exercice\textunderscore 1\textunderscore line\textunderscore search/data/}}, and experiment the learning rate in order to observe convergence or divergence.}

Template files in the folder:
\begin{itemize}
	\item \textbf{main.py}
	\item \textbf{algorithms.py} (to edit)
\end{itemize}

You should observe something like figure \ref{fig:vanilla_gd} in the cases of convergence.

\begin{figure}[htpb]
	\centering
	\includegraphics[width=1\textwidth]{Gradient_descent}
	\caption{Constant step size gradient descent and for several learning rates.}
	\label{fig:vanilla_gd}
\end{figure}

\subsection{Line search: a nested optimization problem}

Considering a fixed iteration step $ \theta_{t}$, we note 

\begin{equation}
    \alpha( \gamma) = \theta_{t}- \gamma \nabla_{\theta}f(\theta_t)\in \mathbb{R}^d
\end{equation}

The \textbf{{exact line seach}}  method attempts to find the optimal step $ \gamma^*$, at each iteration. This means, given the position $ \theta_{t}$, the parameter $ \gamma$ that minimizes the function defined by

\begin{equation}
    \begin{aligned}
        \label{eq:}
	g( \gamma) &= f(\theta_{t}- \gamma \nabla_{\theta}f(\theta_t))\\
	&= f( \alpha( \gamma))
    \end{aligned}
\end{equation}

{\color{Purple}{Is $g: \mathbb{R} \rightarrow \mathbb{R}$ a convex function ?}}
\\

{\color{Purple}{Find the value $\gamma^*$ that minimizes $\gamma\rightarrow g(\gamma)$ for a given
$\theta_t$, and perform gradient descent where $\gamma$ is set optimally thanks to this
method (exact line search). Compare the convergence speeds by measuring the
distance between the iterate and the OLS estimator $\hat{\theta}$ at each iteration, and experiment with the dimensions of the problem. }} 
\\

You can uncomment the lines that call the line\textunderscore search() function in the main file.

You should observe something like figure \ref{fig:ls_gd}

\begin{figure}[htpb]
\hspace*{-2cm}                                                           
	\centering
	\includegraphics[width=1.2\textwidth]{Gradient_descent_n_200_d_100_gaussian}
	\caption{Comparison between constant step size gradient descent and line
    search. In this image, we compare the number of iterations and running time
    to reach a value for $||\theta-\hat{\theta}||$ where $\hat{\theta}$ is the
    OLS estimator and $\theta$ the iterate. The right plot is also in full-log
    scale, as opposed to the previous plots. We note that for these dimensions,
    OLS solved analytically is way faster than GD and line search. For way
    larger dimensions (typically $d>10^5$), the balance might be in favor of
    GD-type methods.}
	\label{fig:ls_gd}
\end{figure}

\begin{figure}[htpb]
\hspace*{-2cm}                                                           
    \centering
    \includegraphics[width=1.2\textwidth]{Gradient_descent_n_200_d_100_gaussian_div}
    \caption{If the learning rate is too large, we might have a divergence !}
    \label{fig:Gradient_descent_n_200_d_100_gaussian_div}
\end{figure}

\begin{figure}[htpb]
\hspace*{-2cm}                                                           
	\centering
	\includegraphics[width=1.2\textwidth]{Gradient_descent_n_1000_d_500_gaussian}
    \caption{Same image for larger dimensions of the problem.}
	\label{fig:ls_gd_500}
\end{figure}





\end{document}
